\documentclass[UTF8,fontset=fandol,zihao=5]{ctexart}

\usepackage[a4paper,margin=2cm]{geometry}
\usepackage{amsmath,amssymb,mathtools}
\usepackage{booktabs,tabularx,array}
\usepackage{xcolor}
\usepackage[most]{tcolorbox}
\usepackage{enumitem}
\usepackage{indentfirst}
\usepackage{hyperref}
\usepackage{tikz}

\usetikzlibrary{
  positioning,
  calc,
  decorations.pathreplacing,
  arrows.meta
}

\usepackage{circuitikz}


% relationbox 中知识关系图的统一样式
\tikzset{
  relnode/.style={
    font=\small,
    align=center
  },
  relline/.style={
    thick
  },
  relarrow/.style={
    -{Stealth},
    thick
  }
}

\DeclareTextFontCommand{\textbf}{\heiti\color{main}}

\definecolor{main}{HTML}{1C551C}
\definecolor{mainlight}{HTML}{FFFFFF}
\definecolor{method}{HTML}{395539}
\definecolor{methodlight}{HTML}{F0FFF0}

\hypersetup{colorlinks=true,linkcolor=main,urlcolor=main}
\setlength{\parindent}{2em}
\linespread{1.5}
\setlength{\abovedisplayskip}{5pt}
\setlength{\belowdisplayskip}{5pt}
\setlength{\abovedisplayshortskip}{3pt}
\setlength{\belowdisplayshortskip}{3pt}
\setlist[itemize]{leftmargin=2em,itemsep=0.1em,topsep=0.2em}
\setlist[enumerate]{leftmargin=2.3em,itemsep=0.1em,topsep=0.2em}

\ctexset{
  section={format=\centering\heiti\zihao{3}\color{main},
           name={第,章},number=\chinese{section},aftername=\quad,
           beforeskip=0.9em,afterskip=0.5em},
  subsection={format=\heiti\zihao{5}\color{main},
              name={,、},number=\chinese{subsection},aftername=,
              beforeskip=0.6em,afterskip=0.22em},
  subsubsection={format=\heiti\zihao{5}\color{main},
                 number={（\chinese{subsubsection}）},aftername=,
                 beforeskip=0.55em,afterskip=0.15em}
}

\newtcolorbox{relationbox}{
  enhanced,breakable,colback=mainlight,colframe=main,
  boxrule=0.65pt,arc=1.4mm,left=2mm,right=2mm,top=1.2mm,bottom=1.2mm
}
\newtcolorbox{methodbox}[1][]{
  enhanced,breakable,colback=methodlight,colframe=method,
  boxrule=0.65pt,arc=1.4mm,left=2mm,right=2mm,top=1.2mm,bottom=1.2mm,
  title=#1,fonttitle=\kaishu,fontupper=\kaishu,
  before upper={\setlength{\parindent}{2em}}
}

\newcommand{\ii}{\mathrm{i}}
\newcommand{\ee}{\mathrm{e}}
\newcommand{\CC}{\mathbb{C}}
\newcommand{\RR}{\mathbb{R}}

% ============================================================
% IEC / GB 逻辑门符号完整封装
%
% AND   : [&]
% OR    : [>=1]
% XOR   : [=1]
% XNOR  : [=]
% NAND  : [&]o
% NOR   : [>=1]o
% NOT   : [1]o
% BUF   : [1]
% AOI22 : 两个 AND 区 + OR 区 + 输出反相圆圈
%
% 使用：
% \ANDgate
% \ORgate
% \XORgate
% \XNORgate
% \NANDgate
% \NORgate
% \NOTgate
% \BUFgate
% \AOIgate
% ============================================================


\newcommand{\logicgateIEC}[1]{%
  \tikz[baseline=-0.6ex]{%
    \node[
      european #1 port,
      scale=0.55,
      transform shape
    ] {};
  }%
}

\newcommand{\logicgateIECinv}[1]{%
  \tikz[baseline=-0.6ex]{%

    % 原始 IEC 门
    \node[
      european #1 port,
      scale=0.55,
      transform shape
    ] (gate) {};

    % 输出反相圆圈
    % 白色填充会覆盖原输出线穿过圆圈的部分
    \filldraw[
      fill=white,
      line width=0.45pt
    ]
      (gate.out)
      circle[radius=0.075em];

    % 圆圈后的输出线
    \draw[line width=0.45pt]
      ([xshift=0.075em]gate.out)
      -- ++(0.28em,0);
  }%
}

% ============================================================
% 可调尺寸
% ============================================================

% 同或门整体尺寸
\newcommand{\IECXNORScale}{1.5}

% AOI22 整体尺寸
\newcommand{\IECAOIScale}{1}


% ============================================================
% 1. 普通二输入 IEC 门
% ============================================================

\newcommand{\IECgatebox}[1]{%
  \tikz[baseline=-0.55ex, line width=0.45pt]{%

    % 输入
    \draw (-0.78em, 0.20em) -- (-0.32em, 0.20em);
    \draw (-0.78em,-0.20em) -- (-0.32em,-0.20em);

    % 主框
    \draw (-0.32em,-0.46em)
          rectangle
          (0.52em,0.46em);

    % 功能符号
    \node[inner sep=0pt]
      at (0.10em,0)
      {\scriptsize #1};

    % 输出
    \draw (0.52em,0) -- (0.98em,0);
  }%
}


% ============================================================
% 2. 带输出反相圆圈的二输入 IEC 门
% ============================================================

\newcommand{\IECgateboxinv}[1]{%
  \tikz[baseline=-0.55ex, line width=0.45pt]{%

    % 输入
    \draw (-0.78em, 0.20em) -- (-0.32em, 0.20em);
    \draw (-0.78em,-0.20em) -- (-0.32em,-0.20em);

    % 主框
    \draw (-0.32em,-0.46em)
          rectangle
          (0.52em,0.46em);

    % 功能符号
    \node[inner sep=0pt]
      at (0.10em,0)
      {\scriptsize #1};

    % 反相圆圈
    \draw (0.64em,0)
          circle[radius=0.12em];

    % 输出
    \draw (0.76em,0) -- (1.08em,0);
  }%
}


% ============================================================
% 3. 基本逻辑门
% ============================================================

% AND
% AND
\newcommand{\ANDgate}{%
  \logicgateIEC{and}%
}

% OR
\newcommand{\ORgate}{%
  \logicgateIEC{or}%
}

% XOR
\newcommand{\XORgate}{%
  \logicgateIEC{xor}%
}

% NAND = AND + 输出反相
\newcommand{\NANDgate}{%
  \logicgateIECinv{and}%
}

% NOR = OR + 输出反相
\newcommand{\NORgate}{%
  \logicgateIECinv{or}%
}


% ============================================================
% 4. BUFFER
% IEC: [1]
% ============================================================

\newcommand{\BUFgate}{%
  \tikz[baseline=-0.55ex, line width=0.45pt]{%

    % 输入
    \draw (-0.68em,0) -- (-0.30em,0);

    % 主框
    \draw (-0.30em,-0.46em)
          rectangle
          (0.50em,0.46em);

    % 功能符号
    \node[inner sep=0pt]
      at (0.10em,0)
      {\scriptsize $1$};

    % 输出
    \draw (0.50em,0) -- (0.88em,0);
  }%
}


% ============================================================
% 5. NOT
% IEC: [1]o
% ============================================================

\newcommand{\IECNOTScale}{1.5}

\newcommand{\NOTgate}{%
  \tikz[
    baseline=-0.55ex,
    scale=\IECNOTScale,
    transform shape,
    line width=0.45pt
  ]{%
    \draw (-0.68em,0) -- (-0.30em,0);

    \draw
      (-0.30em,-0.46em)
      rectangle
      (0.50em,0.46em);

    \node[inner sep=0pt]
      at (0.10em,0)
      {\scriptsize $1$};

    \draw[fill=white]
      (0.62em,0)
      circle[radius=0.12em];

    \draw
      (0.74em,0)
      -- (1.02em,0);
  }%
}


% ============================================================
% 6. XNOR TikZ Shape
%
% IEC: [=]
%
% 可单独作为节点使用：
%
% \node[iec xnor port] (G) {};
%
% 锚点：
%   G.in 1
%   G.in 2
%   G.out
% ============================================================

\makeatletter

\pgfdeclareshape{iec xnor}{

  % --------------------------------------------------------
  % 基本锚点
  % --------------------------------------------------------

  \anchor{center}{
    \pgfpointorigin
  }

  \anchor{text}{
    \pgfpointorigin
  }

  \anchor{north}{
    \pgfpoint{0pt}{0.46em}
  }

  \anchor{south}{
    \pgfpoint{0pt}{-0.46em}
  }

  \anchor{west}{
    \pgfpoint{-0.78em}{0pt}
  }

  \anchor{east}{
    \pgfpoint{0.98em}{0pt}
  }


  % --------------------------------------------------------
  % 输入输出
  % --------------------------------------------------------

  \anchor{in 1}{
    \pgfpoint{-0.78em}{0.20em}
  }

  \anchor{in 2}{
    \pgfpoint{-0.78em}{-0.20em}
  }

  \anchor{out}{
    \pgfpoint{0.98em}{0pt}
  }


  % --------------------------------------------------------
  % 自动边界
  % --------------------------------------------------------

  \anchorborder{%
    \pgf@xa=\pgf@x
    \pgf@ya=\pgf@y
    \pgfpointborderrectangle
      {\pgfpoint{\pgf@xa}{\pgf@ya}}
      {\pgfpoint{0.98em}{0.46em}}
  }


  % --------------------------------------------------------
  % 图形
  % --------------------------------------------------------

  \backgroundpath{

    % 输入 1
    \pgfpathmoveto{
      \pgfpoint{-0.78em}{0.20em}
    }
    \pgfpathlineto{
      \pgfpoint{-0.32em}{0.20em}
    }

    % 输入 2
    \pgfpathmoveto{
      \pgfpoint{-0.78em}{-0.20em}
    }
    \pgfpathlineto{
      \pgfpoint{-0.32em}{-0.20em}
    }

    % 主框
    \pgfpathrectanglecorners
      {\pgfpoint{-0.32em}{-0.46em}}
      {\pgfpoint{0.52em}{0.46em}}

    % 输出
    \pgfpathmoveto{
      \pgfpoint{0.52em}{0pt}
    }
    \pgfpathlineto{
      \pgfpoint{0.98em}{0pt}
    }
  }


  % --------------------------------------------------------
  % 内部符号
  % --------------------------------------------------------

  \foregroundpath{
    \pgftext[
      at=\pgfpoint{0.10em}{0pt}
    ]{
      \normalsize $=$
    }
  }
}

\makeatother


\tikzset{
  iec xnor port/.style={
    shape=iec xnor,
    draw,
    fill=none,
    inner sep=0pt,
    outer sep=0pt,
    line width=0.45pt
  }
}


% 行内 XNOR
\newcommand{\XNORgate}{%
  \tikz[baseline=-0.55ex]{%
    \node[
      iec xnor port,
      scale=\IECXNORScale,
      transform shape
    ] {};
  }%
}


% ============================================================
% 7. AOI22 TikZ Shape
%
% F = \overline{AB + CD}
%
% IEC 结构：
%
%       ┌───┬──────┐
% A ────┤ & │      │
% B ────┤   │ >=1  ├o──
%       ├───┤      │
% C ────┤ & │      │
% D ────┤   │      │
%       └───┴──────┘
%
% 可单独使用：
%
% \node[iec aoi22 port] (G) {};
%
% 锚点：
%   G.in 1
%   G.in 2
%   G.in 3
%   G.in 4
%   G.A
%   G.B
%   G.C
%   G.D
%   G.out
%   G.F
% ============================================================

\makeatletter

\pgfdeclareshape{iec aoi22}{

  % --------------------------------------------------------
  % 基本锚点
  % --------------------------------------------------------

  \anchor{center}{
    \pgfpointorigin
  }

  \anchor{text}{
    \pgfpointorigin
  }

  \anchor{north}{
    \pgfpoint{0pt}{1.25em}
  }

  \anchor{south}{
    \pgfpoint{0pt}{-1.25em}
  }

  \anchor{west}{
    \pgfpoint{-1.65em}{0pt}
  }

  \anchor{east}{
    \pgfpoint{1.65em}{0pt}
  }

  \anchor{north west}{
    \pgfpoint{-1.65em}{1.25em}
  }

  \anchor{north east}{
    \pgfpoint{1.65em}{1.25em}
  }

  \anchor{south west}{
    \pgfpoint{-1.65em}{-1.25em}
  }

  \anchor{south east}{
    \pgfpoint{1.65em}{-1.25em}
  }


  % --------------------------------------------------------
  % 输入锚点
  % --------------------------------------------------------

  \anchor{in 1}{
    \pgfpoint{-1.65em}{0.90em}
  }

  \anchor{in 2}{
    \pgfpoint{-1.65em}{0.50em}
  }

  \anchor{in 3}{
    \pgfpoint{-1.65em}{-0.50em}
  }

  \anchor{in 4}{
    \pgfpoint{-1.65em}{-0.90em}
  }


  % 字母别名

  \anchor{A}{
    \pgfpoint{-1.65em}{0.90em}
  }

  \anchor{B}{
    \pgfpoint{-1.65em}{0.50em}
  }

  \anchor{C}{
    \pgfpoint{-1.65em}{-0.50em}
  }

  \anchor{D}{
    \pgfpoint{-1.65em}{-0.90em}
  }


  % --------------------------------------------------------
  % 输出锚点
  % --------------------------------------------------------

  \anchor{out}{
    \pgfpoint{1.65em}{0pt}
  }

  \anchor{F}{
    \pgfpoint{1.65em}{0pt}
  }


  % --------------------------------------------------------
  % 自动边界
  % --------------------------------------------------------

  \anchorborder{%
    \pgf@xa=\pgf@x
    \pgf@ya=\pgf@y
    \pgfpointborderrectangle
      {\pgfpoint{\pgf@xa}{\pgf@ya}}
      {\pgfpoint{1.65em}{1.25em}}
  }


  % --------------------------------------------------------
  % 图形
  % --------------------------------------------------------

  \backgroundpath{

    % ===== 外框 =====

    \pgfpathrectanglecorners
      {\pgfpoint{-1.20em}{-1.25em}}
      {\pgfpoint{1.15em}{1.25em}}


    % ===== AND / OR 分隔线 =====

    \pgfpathmoveto{
      \pgfpoint{-0.15em}{-1.25em}
    }
    \pgfpathlineto{
      \pgfpoint{-0.15em}{1.25em}
    }


    % ===== 两个 AND 区分隔线 =====

    \pgfpathmoveto{
      \pgfpoint{-1.20em}{0pt}
    }
    \pgfpathlineto{
      \pgfpoint{-0.15em}{0pt}
    }


    % ===== 输入引线 =====

    \pgfpathmoveto{
      \pgfpoint{-1.65em}{0.90em}
    }
    \pgfpathlineto{
      \pgfpoint{-1.20em}{0.90em}
    }

    \pgfpathmoveto{
      \pgfpoint{-1.65em}{0.50em}
    }
    \pgfpathlineto{
      \pgfpoint{-1.20em}{0.50em}
    }

    \pgfpathmoveto{
      \pgfpoint{-1.65em}{-0.50em}
    }
    \pgfpathlineto{
      \pgfpoint{-1.20em}{-0.50em}
    }

    \pgfpathmoveto{
      \pgfpoint{-1.65em}{-0.90em}
    }
    \pgfpathlineto{
      \pgfpoint{-1.20em}{-0.90em}
    }


    % ===== 输出反相圆圈 =====

    \pgfpathmoveto{
      \pgfpoint{1.15em}{0pt}
    }
    \pgfpathlineto{
      \pgfpoint{1.19em}{0pt}
    }

    \pgfpathcircle
      {\pgfpoint{1.30em}{0pt}}
      {0.11em}

    \pgfpathmoveto{
      \pgfpoint{1.41em}{0pt}
    }
    \pgfpathlineto{
      \pgfpoint{1.65em}{0pt}
    }
  }


  % --------------------------------------------------------
  % 内部符号
  % --------------------------------------------------------

  \foregroundpath{

    % 上 AND
    \pgftext[
      at=\pgfpoint{-0.68em}{0.63em}
    ]{
      \scriptsize \&
    }

    % 下 AND
    \pgftext[
      at=\pgfpoint{-0.68em}{-0.63em}
    ]{
      \scriptsize \&
    }

    % OR
    \pgftext[
      at=\pgfpoint{0.50em}{0.53em}
    ]{
      \scriptsize $\geq 1$
    }
  }
}

\makeatother


\tikzset{
  iec aoi22 port/.style={
    shape=iec aoi22,
    draw,
    fill=none,
    inner sep=0pt,
    outer sep=0pt,
    line width=0.45pt
  }
}


% 行内 AOI22
\newcommand{\AOIgate}{%
  \tikz[baseline=-0.55ex]{%
    \node[
      iec aoi22 port,
      scale=\IECAOIScale,
      transform shape
    ] {};
  }%
}

% =========================================================
% IEEE 基础逻辑门封装
% =========================================================

\newcommand{\logicgateIEEE}[1]{%
  \tikz[baseline=-0.6ex]{%
    \node[
      ieeestd #1 port,
      scale=0.55
    ] {};
  }%
}
% 与门
\newcommand{\ANDEgate}{%
  \logicgateIEEE{and}%
}

% 或门
\newcommand{\OREgate}{%
  \logicgateIEEE{or}%
}

% 异或门
\newcommand{\XOREgate}{%
  \logicgateIEEE{xor}%
}

% 与非门
\newcommand{\NANDEgate}{%
  \logicgateIEEE{nand}%
}

% 或非门
\newcommand{\NOREgate}{%
  \logicgateIEEE{nor}%
}

% 同或门
\newcommand{\XNOREgate}{%
  \logicgateIEEE{xnor}%
}

% 缓冲器
\newcommand{\BUFEgate}{%
  \logicgateIEEE{buffer}%
}

% 非门
\newcommand{\NOTEgate}{%
  \logicgateIEEE{not}%
}

\newcommand{\AOIEgate}{%
  \tikz[baseline=-0.6ex, scale=0.38, transform shape]{%

    % 两个与门
    \node[ieeestd and port] (and1) at (0, 0.72) {};
    \node[ieeestd and port] (and2) at (0,-0.72) {};

    % 或非门
    \node[ieeestd nor port] (nor) at (2.8,0) {};

    % 与门 → 或非门
    \draw
      (and1.out) -- ++(0.55,0) |- (nor.in 1)
      (and2.out) -- ++(0.55,0) |- (nor.in 2);

  }%
}

\begin{document}

\begin{center}
  {\heiti\zihao{1}\color{main} 数字电路笔记}
\end{center}

\vspace{0.45em}

\begin{relationbox}
\[
  \boxed{\text{预备知识}}
  \longrightarrow
  \boxed{\text{组合逻辑电路}}
  \longrightarrow
  \boxed{\text{时序逻辑电路}}
\]
\end{relationbox}

不连续、是固定量整数倍的量称为数字量，与数字量对应的信号称为数字信号，处理数字信号的电路称为数字电路。

\section{预备知识}

\subsection{数制}

\begin{enumerate}
  \item 不同的数制
  \begin{center}
    \begin{tabularx}{0.92\textwidth}{>{\bfseries}p{2.5cm}X X}
      \toprule
      数制（符号） & 基数 & 数码 \\
      \midrule
      二进制 B & 2 & 0,1 \\
      八进制 O & 8 & 0-7 \\
      十进制 D & 10 & 0-9 \\
      十六进制 H & 16 & 0-9, A-F \\
      \bottomrule
    \end{tabularx}
  \end{center}

  \item 数制转换
\begin{methodbox}[非十进制互转]
每3位或4位分组，然后展开或缩写即可。例：
\[
\left ( 11001101.011 \right )_2 = \left ( 315.3 \right )_8 = \left ( CD.6 \right )_{16}
\]
\end{methodbox}

\begin{methodbox}[非十进制转十进制]
按权展开即可。例：
\[
\left ( 101.01 \right )_2 = 1 \times 2^2 + 0 \times 2^1 + 1 \times 2^0 + 0 \times 2^{-1} + 1 \times 2^{-2} = 5.25
\]
\[
\left ( 73.4 \right )_8 = 7 \times 8^1 + 3 \times 8^0 + 4 \times 8^0 + 4 \times 8^{-1} = 59.5
\]
\[
\left ( 3\mathrm{D}.\mathrm{C} \right )_{16} = 3 \times 16^1 + 13 \times 16^0 + 12 \times 16^{-1} = 61.75
\]
\end{methodbox}

\begin{methodbox}[十进制转非十进制]
整数部分除基取余，从低到高；小数部分乘基取整，从高到低。例：13.25

13除以2余1，6除以2余0，3除以2余1，1除以2余1，整数部分为1101；

0.25乘以2得0.5，取整为0，0.5乘以2得1.0，取整为1，小数部分为01；

故$13.25 = \left ( 1101.01 \right )_2$
\end{methodbox}
\end{enumerate}

\subsection{编码}

\begin{enumerate}
  \item 原码、反码与补码
  \begin{center}
    \begin{tabularx}{0.92\textwidth}{>{\bfseries}p{2.5cm}X X}
      \toprule
      编码类型 & 描述 & 缺点 \\
      \midrule
      原码 & 首位表示符号位，剩余位为二进制绝对值 & 减法运算困难 \\
      反码 & 正数与零不变；负数的符号位不变，其余位为原码取反 & 有两个反码表示零 \\
      补码 & 正数与零不变；负数的补码为反码加1 & 无 \\
      \bottomrule
    \end{tabularx}
  \end{center}

  \item 二-十进制编码（BCD码）
  \begin{center}
    \begin{tabularx}{0.92\textwidth}{>{\bfseries}p{2.5cm}X X}
      \toprule
      编码类型 & 权 & 特性 \\
      \midrule
      8421码 & 8, 4, 2, 1 & 简便直观 \\
      2421码 & 2, 4, 2, 1 & 对9自补码 \\
      余3码 & 无 & 对9自补码 \\
      5421码 & 5, 4, 2, 1 & 可以产生对称方波输出 \\
      Gray码 & 无 & 相邻两数码只有一位二进制数不同 \\
      \bottomrule
    \end{tabularx}
  \end{center}

  \item 可靠性编码
\begin{enumerate}[label=(\arabic*)]
  \item Gray码

Gray码是一种特殊的二进制编码，其特点是相邻两个数码之间只有一位二进制数不同。
Gray码的主要优点是减少了在数字电路中由于信号变化引起的误差。

4位反射Gray码的编码表如下：
\begin{table}[htbp]
    \centering
    \begin{tabular}{|c|c|c|c|}
        \hline
        $0\;(0000)$  & $1\;(0001)$  & $2\;(0011)$  & $3\;(0010)$  \\
        \hline
        $4\;(0110)$  & $5\;(0111)$  & $6\;(0101)$  & $7\;(0100)$  \\
        \hline
        $8\;(1100)$  & $9\;(1101)$  & $10\;(1111)$ & $11\;(1110)$ \\
        \hline
        $12\;(1010)$ & $13\;(1011)$ & $14\;(1001)$ & $15\;(1000)$ \\
        \hline
    \end{tabular}
\end{table}

  \item 奇偶校验码

在要传输的数据中增加一个校验位，使得数据中1的个数为奇数或偶数，从而实现对数据传输错误的检测。

最常用的是奇校验码，因其使得全0成为非法码，便于检测。

  \item 其他可靠性编码：Hamming码等

\end{enumerate}
\end{enumerate}

\subsection{逻辑代数}

逻辑常量：0, 1 \\
基本逻辑运算：与（AND）、或（OR）、非（NOT）

\subsubsection{逻辑运算}

\begin{enumerate}

  \item 基本逻辑运算

  \begin{center}
  \renewcommand{\arraystretch}{1.4}
  \begin{tabularx}{0.92\textwidth}{
    >{\bfseries}c
    >{\centering\arraybackslash}X
    >{\centering\arraybackslash}X
    >{\centering\arraybackslash}p{2.5cm}
    >{\centering\arraybackslash}p{2.5cm}
  }
    \toprule
    逻辑门 & 运算规则 & 逻辑运算式 & IEEE图形符号 & IEC图形符号 \\
    \midrule
    与门（AND） & 有0则0 & $F \left ( A , \, B \right ) = A \cdot B$ & \ANDEgate & \ANDgate \\
    或门（OR） & 有1则1 & $F \left ( A , \, B \right ) = A + B$ & \OREgate & \ORgate \\
    非门（NOT） & 反转输入 & $F \left ( A \right ) = \overline{A}$ & \NOTEgate & \NOTgate \\
    \bottomrule
  \end{tabularx}
\end{center}

\item 复合逻辑运算

  \begin{center}
  \renewcommand{\arraystretch}{1.4}
  \begin{tabularx}{0.92\textwidth}{
    >{\bfseries}c
    >{\centering\arraybackslash}X
    >{\centering\arraybackslash}X
    >{\centering\arraybackslash}p{2.5cm}
    >{\centering\arraybackslash}p{2.5cm}
  }
    \toprule
    逻辑门 & 运算规则 & 逻辑运算式 & IEEE图形符号 & IEC图形符号 \\
    \midrule
    与非门（NAND） & 有0则1 & $F \left ( A , \, B \right ) = \overline{A \cdot B}$ & \NANDEgate & \NANDgate \\
    或非门（NOR） & 有1则0 & $F \left ( A , \, B \right ) = \overline{A + B}$ & \NOREgate & \NORgate \\
    与或非门（AOI） & 与、或、非 & $F \left ( A , \, B , \, C , \, D \right ) = \overline{AB + CD}$ & \AOIEgate & \AOIgate \\
    异或门（XOR） & 同0异1 & $F \left ( A , \, B \right ) = A \oplus B = \overline{A} \cdot B + A \cdot \overline{B}$ & \XOREgate & \XORgate \\
    同或门（XNOR） & 同1异0 & $F \left ( A , \, B \right ) = A \odot B = \overline{A} \cdot \overline{B} + A \cdot B$ & \XNOREgate & \XNORgate \\
    \bottomrule
  \end{tabularx}
\end{center}

\end{enumerate}

\subsubsection{逻辑代数公式}

\begin{enumerate}

\item 基本公式

  \begin{center}
    \begin{tabularx}{0.92\textwidth}{>{\bfseries}p{2.5cm}X X}
      \toprule
      名称 & 或式 & 与式 \\
      \midrule
      交换律 & $A + B = B + A$ & $A \cdot B = B \cdot A$ \\
      结合律 & $(A + B) + C = A + (B + C)$ & $(A \cdot B) \cdot C = A \cdot (B \cdot C)$ \\
      分配律 & $A + (B \cdot C) = (A + B) \cdot (A + C)$ & $A \cdot (B + C) = (A \cdot B) + (A \cdot C)$ \\
      吸收律 & $A + (A \cdot B) = A$ & $A \cdot (A + B) = A$ \\
      0-1律 & $A + 1 = 1, \qquad A + 0 = A$ & $A \cdot 0 = 0, \qquad A \cdot 1 = A$ \\
      互补律 & $A + \overline{A} = 1$ & $A \cdot \overline{A} = 0$ \\
      重叠律 & $A + A = A$ & $A \cdot A = A$ \\
      还原律 & $A + \overline{A} = 1$ & $A \cdot \overline{A} = 0$ \\
      反演律 & $\overline{A + B} = \overline{A} \cdot \overline{B}$ & $\overline{A \cdot B} = \overline{A} + \overline{B}$ \\
      \bottomrule
    \end{tabularx}
  \end{center}

  \item 常用公式

    \begin{itemize}
      \item $A \cdot B + A \cdot \overline{B} = A$
      \item $A + \overline{A} \cdot B = A + B$
      \item $A \cdot B + \overline{A} \cdot C + B \cdot C = A \cdot B + \overline{A} \cdot C$ 
      \item $ \left ( A + B \right ) \cdot \left (\overline{A} + C \right ) = A \cdot C + \overline{A} \cdot B$
      \item $ \overline{\overline{A} \cdot B + A \cdot \overline{B}} = \overline{A} \cdot \overline{B} + A \cdot B$，
      即$\overline{A \oplus B} = A \odot B$\\
    \end{itemize}

\end{enumerate}

\subsubsection{逻辑代数基本规则}

    \begin{enumerate}
      \item 代入规则：用逻辑表达式$F$代入某个变量$X$，新逻辑表达式仍然成立。
      \item 反演规则

      将原函数$F$作如下代换：

      \begin{itemize}
        \item $ + \leftrightarrow \cdot $
        \item $ 1 \leftrightarrow 0 $
        \item $ A \leftrightarrow \overline{A} $
      \end{itemize}

      即得到反函数$\overline{F}$。

      \item 对偶规则

      将原函数$F$作如下代换：

      \begin{itemize}
        \item $ + \leftrightarrow \cdot $
        \item $ 1 \leftrightarrow 0 $
      \end{itemize}

      即得到对偶函数$F'$。

    \end{enumerate}

\subsubsection{正负逻辑}

    \begin{itemize}
      \item 正逻辑：逻辑1对应高电平，逻辑0对应低电平。
      \item 负逻辑：逻辑1对应低电平，逻辑0对应高电平。
    \end{itemize}

    \begin{methodbox}[正负逻辑转换]
      对自变量预先取反，即能将负逻辑电路转换为正逻辑电路，反之亦然。
    \end{methodbox}

\subsubsection{逻辑函数}

\end{document}